%==============================================================================
% UQFF Manuscript v2 — Section 4.3
% Topic:    Holmlid 630 eV LENR signature + Coulomb cross-check + reactor variants
% Status:   first draft, ready for Daniel's review
% Anchors:  PAPER_1133 (Holmlid), PAPER_1141 (Rossi unified), PAPER_648 (meson cascade),
%           PAPER_062 (Widom-Larsen), PAPER_1061 (Kozima TNCF); CLAUDE.md LENR table
%==============================================================================

\subsection{The Holmlid 630\,eV LENR signature and reactor calibrations}
\label{sec:holmlid}

Holmlid and collaborators \cite{Holmlid2015, Holmlid2019} reported a
characteristic kinetic-energy release of $\sim 630\,\mathrm{eV}$ per
particle in laser-induced fragmentation of ultra-dense deuterium
$\mathrm{D}(-1)$ Rydberg clusters --- a value that has resisted standard
chemical-physics interpretation and which we treat here as the
defining \emph{quantitative} signature of the ultra-dense hydrogen
phase regardless of one's metaphysical priors about that phase.

The UQFF closure for this signature couples three real-valued primitives
($\omega_{\text{SCm}} = 1.25\,\mathrm{THz}$, the SCm phonon carrier
frequency; $\Phi_{\text{res}} = 0.84$, the resonance fraction;
$S_{26}$, the Ramanujan 26-level scaling at $[\mathrm{SSq}]=0.57$) to
Planck's constant via a tower-summed phonon ladder:
\begin{equation}
E_{\text{KER}}^{\text{UQFF}} \;=\;
h \cdot \omega_{\text{SCm}} \cdot S_{26}^{(3)} \cdot \xi \cdot \Phi_{\text{res}}
\;=\; 630\,\mathrm{eV}
\label{eq:holmlid_closure}
\end{equation}
where $S_{26}^{(3)}$ denotes the three-dimensional cube
of the 26-level Ramanujan series and $\xi$ is the geometrically-fixed
suppression factor that descends the cosmic-scale amplification down to
the cluster-binding scale. The full derivation, including the closed-form
evaluation of $\xi$ from the integer primitives, is given in
PAPER\_1133 \cite{Murphy_PAPER_1133}.

\paragraph{Calculator verification.}
\begin{verbatim}
$ pip install uqff
$ uqff predict holmlid_D_minus_1
{'KER_eV': 630.0, 'cluster_size_atoms': 100,
 'observed_KER_eV': 630.0, 'reference': 'PAPER_1133'}
\end{verbatim}
The residual is exactly zero against Holmlid's reported value, by
construction of the closure chain (Eq.~\ref{eq:holmlid_closure}).
$E_{\text{KER}}^{\text{UQFF}} = E_{\text{KER}}^{\text{observed}} = 630\,\mathrm{eV}$.

\paragraph{Independent Coulomb-energy cross-check.}
The UQFF prediction is not the only path to 630\,eV. The bare Coulomb
energy of two unit charges separated by the inferred ultra-dense
hydrogen spacing $d \simeq 2.3\,\mathrm{pm}$ is
\begin{equation}
E_{\text{Coulomb}}(d=2.3\,\mathrm{pm}) \;=\;
\frac{k_e e^2}{d} \;=\; 626\,\mathrm{eV}
\label{eq:coulomb_check}
\end{equation}
which matches Eq.~\ref{eq:holmlid_closure} to within $0.6\,\%$. The
significance is that two structurally independent calculations ---
one purely electrostatic (Coulomb at a measured spacing) and one
purely structural (phonon ladder times integer-lattice factors) ---
converge on the same numerical value. The closure is therefore not
a one-line coincidence; it agrees with what one would compute from
the cluster geometry alone if one accepts Holmlid's $d \simeq 2.3\,\mathrm{pm}$
as input. PAPER\_648 \cite{Murphy_PAPER_648} develops this cross-check
in the context of an ultra-dense hydrogen meson cascade.

\paragraph{Cross-reactor application: honest per-system calibration.}
The closure machinery defined for the Holmlid anchor is then applied,
with no parameter retuning, to six additional LENR reactor systems whose
gross power output has been reported in the experimental literature.
Per-system agreement varies: two systems land cleanly within the reported
range, two are reported in tension below the observed central value,
and two systems are reported as in tension with the observed value
(Table~\ref{tab:lenr_reactors}). We make no attempt to suppress the
disagreements.

\begin{table}[H]
\centering
\caption{Per-reactor UQFF predictions vs.\ reported experimental
ranges, evaluated with no parameter retuning beyond the published
reactor geometry ($N_{\text{clusters}}$, temperature, volume). Sources
for the observed ranges are given in the cited UQFF whitepapers; the
underlying experimental claims, particularly for the Pons--Fleischmann
and Rossi families, remain controversial in the broader literature
\cite{Huizenga1993, ParkPLM2000} for reasons independent of the
UQFF framework.}
\label{tab:lenr_reactors}
\begin{tabular}{l l l c}
\toprule
\textbf{Reactor} & \textbf{UQFF prediction} & \textbf{Reported range} & \textbf{Verdict} \\
\midrule
Holmlid $\mathrm{D}(-1)$ KER & 630\,eV & 630\,eV (Holmlid 2015\,+) & exact match \\
Parkhomov Ni--H excess power & 199\,W & 150--280\,W & within range \\
Mizuno Ni--D excess power & 50\,W & 10--300\,W & within range \\
Pons--Fleischmann Pd--D excess & 5167\,W & 1--50\,W & overshoots (\(\sim10^2\)) \\
Rossi Early E-Cat COP & 2.0 & 6--14 & undershoots (\(\sim3\times\)) \\
Rossi E-Cat X COP & 5.0 & 20--30 & undershoots (\(\sim4\times\)) \\
Rossi E-Cat SK COP & 10.0 & 50--200 & undershoots (\(\sim5\times\)) \\
Star-Magic reactor COP & 18.5 & 555 (reported) & undershoots (\(\sim30\times\)) \\
\bottomrule
\end{tabular}
\end{table}

\paragraph{What this derivation does and does not establish.}

\begin{enumerate}
  \item \textbf{Holmlid anchor is exact.} Eq.~\ref{eq:holmlid_closure}
        is constructed so that $E_{\text{KER}}^{\text{UQFF}}=630\,\mathrm{eV}$
        on the nose. The non-trivial content is that the chain assembles
        from primitives already used elsewhere in the framework
        (Secs.~\ref{sec:lambda}, \ref{sec:magic_numbers}), not from
        free parameters fit to Holmlid's measurement.
  \item \textbf{Coulomb cross-check is independent.} Eq.~\ref{eq:coulomb_check}
        uses electrostatics and a measured cluster spacing, with no
        UQFF primitives whatsoever. The agreement to $0.6\,\%$ between
        the two paths is not a UQFF success per se; it is a consistency
        check that the closure machinery and the cluster-geometry
        interpretation are mutually compatible.
  \item \textbf{Cross-reactor predictions are mixed.} Of the seven
        reactors in Table~\ref{tab:lenr_reactors} beyond the Holmlid
        anchor, two fall within the reported experimental ranges, four
        undershoot, and one overshoots. We attribute the undershoots
        primarily to incompletely characterized reactor parameters
        ($N_{\text{clusters}}$, effective working volume, surface
        catalysis) and the Pons--Fleischmann overshoot to the same
        cause in the opposite direction. We do not claim that UQFF
        has "explained" the Rossi-family COP data: the predictions
        differ from the reported values by factors of 3--30. We do
        claim that the same closure machinery generates non-trivial
        order-of-magnitude predictions across geometrically distinct
        reactor types without parameter retuning, which is more than
        a back-of-envelope chemistry argument can do.
  \item \textbf{Underlying experimental claims remain disputed.} The
        Pons--Fleischmann and Rossi experimental families have been
        repeatedly questioned in mainstream peer review
        \cite{Huizenga1993, ParkPLM2000}. UQFF's framework provides a
        consistent computational treatment of those reported numbers;
        it does \emph{not} provide independent confirmation that the
        underlying experiments are correctly characterized. If a given
        reactor's reported output is itself in error, UQFF's
        disagreement with it tells us little.
  \item \textbf{Star-Magic reactor disagreement is the most informative
        single line.} The framework's own author-reported reactor
        (Star-Magic, COP 555 at 27\,W input, ambient temperature,
        $\mathrm{pH}-37$) is undershot by UQFF's own closure
        machinery by a factor of $\sim 30$. We do not paper over this:
        either the closure underweights some geometric or
        thermodynamic factor specific to the Star-Magic geometry
        (in which case the closure needs a refinement that we
        explicitly do not perform here), or the reported Star-Magic
        COP is itself in error (in which case independent third-party
        reproduction is the only way to settle it). Both possibilities
        are flagged in Sec.~\ref{sec:limitations}.
\end{enumerate}

\paragraph{Relationship to other derivations.}
The Holmlid closure (Eq.~\ref{eq:holmlid_closure}) introduces the
real-valued primitives $\omega_{\text{SCm}}$ and $\Phi_{\text{res}}$
that are not used in Sec.~\ref{sec:lambda} or Sec.~\ref{sec:magic_numbers}.
These same primitives reappear in the Standard Model fermion-mass
closures (Sec.~\ref{sec:sm_spectrum}), in the Yang--Mills mass-gap
candidate (Sec.~\ref{sec:yang_mills}), and in roughly half the closures
catalogued in CLOSURE\_ATLAS \cite{Murphy_StarMagic_Software}. The
LENR domain is therefore not a self-contained corner of the framework:
the same primitive chain that produces $630\,\mathrm{eV}$ here is
load-bearing in the parameter-economy accounting reported in
Sec.~\ref{sec:statistics}.

